only if "@@M0@@ is torsion-free", i.e. @@M1@@ is isomorphic to a product of groups @@M2@@ and @@M3@@. Thus in this case, and in this case alone, the homomorphism
@@M4@@
is an isomorphism. Thus, when @@M5@@ has torsion, the underlying scheme provides an example of a nonbijective @@M6@@.

\sgaStableTarget{sga6.e14.para.4-8}\textbf{4.8.} Similarly, one finds an example in which the homomorphism @@M7@@ of (4.1)
\sgaStableTarget{sga6.e14.eq.4-9}@@M8@@
does not factor through @@M9@@, or, equivalently here, an example in which this homomorphism is nonzero in degrees @@M10@@. Restricting ourselves (as suitable specialization theorems allow) to the case @@M11@@, we are reduced to studying the homomorphism
\sgaStableTarget{sga6.e14.eq.4-10}@@M12@@
which is moreover deduced by passage to the quotient from the canonical homomorphism
@@M13@@
It is tempting to conjecture that the homomorphism (4.10) is in fact injective, since the torsion primes of @@M14@@ are known to be exactly those dividing the order of the finite group @@M15@@. Were this so, as soon as @@M16@@ had torsion the homomorphism (4.10) would not factor through @@M17@@, and thus, for suitable @@M18@@ (more precisely, for @@M19@@ dividing the order of @@M20@@, i.e. @@M21@@ a torsion coefficient of @@M22@@), the same would hold for (4.9). One can at least assert that, for the group @@M23@@ over the field @@M24@@, the homomorphism (4.10) is nonzero in degree @@M25@@; more precisely, there is an element of @@M26@@ whose image in @@M27@@ has order @@M28@@; consequently, the homomorphism (4.9) does not factor through @@M29@@ when @@M30@@. The fact used concerning @@M31@@, which amounts to the existence of an element of @@M32@@ whose image in @@M33@@ is nonzero, is found in [3], using the familiar dictionary between the languages of complex semisimple algebraic groups and compact Lie groups; it was kindly communicated to me by J.P. SERRE.

